% Options for packages loaded elsewhere
\PassOptionsToPackage{unicode}{hyperref}
\PassOptionsToPackage{hyphens}{url}
\documentclass[
  9pt,
  ignorenonframetext,
]{beamer}
\newif\ifbibliography
\usepackage{pgfpages}
\setbeamertemplate{caption}[numbered]
\setbeamertemplate{caption label separator}{: }
\setbeamercolor{caption name}{fg=normal text.fg}
\beamertemplatenavigationsymbolsempty
% remove section numbering
\setbeamertemplate{part page}{
  \centering
  \begin{beamercolorbox}[sep=16pt,center]{part title}
    \usebeamerfont{part title}\insertpart\par
  \end{beamercolorbox}
}
\setbeamertemplate{section page}{
  \centering
  \begin{beamercolorbox}[sep=12pt,center]{section title}
    \usebeamerfont{section title}\insertsection\par
  \end{beamercolorbox}
}
\setbeamertemplate{subsection page}{
  \centering
  \begin{beamercolorbox}[sep=8pt,center]{subsection title}
    \usebeamerfont{subsection title}\insertsubsection\par
  \end{beamercolorbox}
}
% Prevent slide breaks in the middle of a paragraph
\widowpenalties 1 10000
\raggedbottom
\AtBeginPart{
  \frame{\partpage}
}
\AtBeginSection{
  \ifbibliography
  \else
    \frame{\sectionpage}
  \fi
}
\AtBeginSubsection{
  \frame{\subsectionpage}
}
\usepackage{iftex}
\ifPDFTeX
  \usepackage[T1]{fontenc}
  \usepackage[utf8]{inputenc}
  \usepackage{textcomp} % provide euro and other symbols
\else % if luatex or xetex
  \usepackage{unicode-math} % this also loads fontspec
  \defaultfontfeatures{Scale=MatchLowercase}
  \defaultfontfeatures[\rmfamily]{Ligatures=TeX,Scale=1}
\fi
\usefonttheme[]{serif}
\usepackage[]{mathpazo}
\ifPDFTeX\else
  % xetex/luatex font selection
\fi
% Use upquote if available, for straight quotes in verbatim environments
\IfFileExists{upquote.sty}{\usepackage{upquote}}{}
\IfFileExists{microtype.sty}{% use microtype if available
  \usepackage[]{microtype}
  \UseMicrotypeSet[protrusion]{basicmath} % disable protrusion for tt fonts
}{}
\makeatletter
\@ifundefined{KOMAClassName}{% if non-KOMA class
  \IfFileExists{parskip.sty}{%
    \usepackage{parskip}
  }{% else
    \setlength{\parindent}{0pt}
    \setlength{\parskip}{6pt plus 2pt minus 1pt}}
}{% if KOMA class
  \KOMAoptions{parskip=half}}
\makeatother
\usepackage{graphicx}
\makeatletter
\newsavebox\pandoc@box
\newcommand*\pandocbounded[1]{% scales image to fit in text height/width
  \sbox\pandoc@box{#1}%
  \Gscale@div\@tempa{\textheight}{\dimexpr\ht\pandoc@box+\dp\pandoc@box\relax}%
  \Gscale@div\@tempb{\linewidth}{\wd\pandoc@box}%
  \ifdim\@tempb\p@<\@tempa\p@\let\@tempa\@tempb\fi% select the smaller of both
  \ifdim\@tempa\p@<\p@\scalebox{\@tempa}{\usebox\pandoc@box}%
  \else\usebox{\pandoc@box}%
  \fi%
}
% Set default figure placement to htbp
\def\fps@figure{htbp}
\makeatother
\setlength{\emergencystretch}{3em} % prevent overfull lines
\providecommand{\tightlist}{%
  \setlength{\itemsep}{0pt}\setlength{\parskip}{0pt}}
% Document Class options:
% - You can use article(beamer) or presentation(beamer)
% - In presentation, you can use following options:
%    'trans' for transparent, 'handout' for printing-friendly version
%    'notes' also output notes, 'notesonly' only output notes

% If you have several lectures defined in the same document, uncomment this:
%\includeonlylecture{cours02}

\usepackage{fontspec}
%\fontfamily{ppl}\selectfont
\setmainfont{Arial}
%\setmainfont{DejaVu LGC Sans}
%\setmainfont[SmallCapsFont=Fontin SmallCaps]{Fontin-Regular}
%\setmainfont{Gentium Basic}
%\setsansfont{DejaVu LGC Sans}
%\setmonofont{DejaVu LGC Sans Mono}
%\setmainfont{Calibri}
%\setmainfont{Cambria}
%\setsansfont{Candara}
%\setmonofont{Consolas}

% UMONS-EcoNum template configuration
\mode<all>
{
  \definecolor{umons-red}{RGB}{168, 0, 57}
  \definecolor{umons-turquoise}{RGB}{0, 171, 204}
  \definecolor{umons-gray}{RGB}{150, 150, 150}
  \definecolor{umons-gray2}{RGB}{110, 110, 110}
  \usepackage[labelfont=bf, tableposition=top]{caption}
}
\mode<article>
{
  \usepackage{beamerbasearticle}
  \usepackage{fullpage}
  \pdfimageresolution 96
}
\mode<beamer>
{
  % Customizable items (make sure corresponding pictures are in /template subdir):
  % The background image
  \usebackgroundtemplate{\includegraphics[width=1.07\paperwidth,
height=1.07\paperheight]{../../template/background-hex.jpg}}
  % The logo picture
  \logo{\includegraphics[width=2cm]{../../template/UMONS-logo.pdf}}
}
\mode<handout>
{
  \definecolor{umons-red}{RGB}{84, 84, 84}
  \definecolor{umons-turquoise}{RGB}{168, 168, 168}
  \usepackage{pgfpages}
  \pgfpagesuselayout{2 on 1}[a4paper, border shrink=5mm]
}
\mode<presentation>
{
  % General theme to use
  \usetheme{CambridgeUS}
  %\usetheme{AnnArbor}
  %\usetheme{default}
  %\usetheme{Pittsburgh}
  %\usetheme{Dresden}
  %\usetheme{Hannover}
  % Color theme (but must colors redefined hereunder)
  \usecolortheme{beaver}
  % Theme used for bullets
  \useinnertheme{rectangles}
  %or \useinnertheme{circles}
  % This is to get semi-transparent covered items
  \setbeamercovered{transparent}
  % Rounded rectangular blocks without shadow
  \setbeamertemplate{blocks}[rounded][shadow=false]
 %\usepackage{beamerthemeshadow}

  \setbeamersize{text margin left=1em,text margin right=1em}
  \setbeamerfont{title}{size=\huge}
  \setbeamerfont{subtitle}{size=\large}
  \setbeamerfont{institute}{size=\small}
  \setbeamerfont{date}{size=\small}
  \setbeamerfont{framesubtitle}{size=\small}
  % This stands out a little bit more than umons-red!
  \setbeamercolor{alerted text}{fg=red!80!black}
  %\setbeamercolor{alerted text}{fg=umons-red}
  \setbeamercolor*{palette primary}{fg=black, bg=umons-turquoise}
  \setbeamercolor*{palette secondary}{fg=black, bg=umons-turquoise}
  \setbeamercolor*{palette tertiary}{fg=umons-gray!30, bg=umons-red!80!black}
  \setbeamercolor*{palette quaternary}{fg=black, bg=umons-turquoise}
  \setbeamercolor*{upper separation line head left}{parent=palette tertiary}
  \setbeamercolor*{upper separation line head right}{parent=palette primary}
  \setbeamercolor{title}{fg=umons-red}
  \setbeamercolor*{titlelike}{fg=umons-red}
  \setbeamercolor{institute}{fg=umons-gray2}
  \setbeamercolor{date}{fg=umons-gray2}
  \setbeamercolor{frametitle}{fg=umons-red, bg=white}
  \setbeamercolor{frametitle right}{bg=white}
  \setbeamercolor{structure}{fg=umons-turquoise}
  \setbeamercolor*{separation line}{}
  \setbeamercolor*{fine separation line}{}
  \setbeamercolor*{sidebar}{fg=umons-red,bg=white}
  \setbeamercolor*{palette sidebar primary}{fg=umons-turquoise!10!black}
  \setbeamercolor*{palette sidebar secondary}{fg=umons-red}
  \setbeamercolor*{palette sidebar tertiary}{fg=umons-red!50!black}
  \setbeamercolor*{palette sidebar quaternary}{fg=umons-red}
  \setbeamercolor*{example text}{fg=umons-gray2}

  % Change the footer
  \makeatother
  \setbeamertemplate{footline}
  {
    \leavevmode%
    \hbox{%
    \begin{beamercolorbox}[wd=.35\paperwidth,ht=2.25ex,dp=1ex,center]{author in head/foot}%
      \usebeamerfont{author in head/foot}\insertshortauthor
    \end{beamercolorbox}%
    \begin{beamercolorbox}[wd=.57\paperwidth,ht=2.25ex,dp=1ex,center]{title in head/foot}%
      \usebeamerfont{title in head/foot}\insertshorttitle\hspace*{3em}
    \end{beamercolorbox}%
    \begin{beamercolorbox}[wd=.08\paperwidth,ht=2.25ex,dp=1ex,center]{title in head/foot}%
      \insertframenumber{} / \inserttotalframenumber\hspace*{1ex}
    \end{beamercolorbox}}%
    \vskip0pt%
  }
  \makeatletter
  \setbeamertemplate{navigation symbols}{}

  % Custom Sweave styles
  %\definecolor{SinputColor}{RGB}{25, 25, 125}
  %\definecolor{SoutputColor}{RGB}{110, 110, 110}
  %\DefineVerbatimEnvironment{Sinput}{Verbatim}{fontsize=\small,formatcom=\color{SinputColor}}
  %\DefineVerbatimEnvironment{Soutput}{Verbatim}{fontsize=\small,formatcom=\color{SoutputColor}}
  %\DefineVerbatimEnvironment{Scode}{Verbatim}{fontsize=\small,\color{SinputColor}}
}

% To get coherent font between text and equations
\usefonttheme{professionalfonts}
%\usepackage{mathpazo}
% Useful package in this context
\usepackage{listings}
% Command used to place a picture somewhere in the slide (relative to current pos)
\newcommand{\putat}[3]{\begin{picture}(0,0)(0,0)\put(#1,#2){#3}\end{picture}}

% Idem, but in absolute pos (slide = 128x96mm)
\usepackage{textpos}
\newcommand{\putabs}[3]{\begin{textblock*}{10mm}[0,0](#1,#2){#3}\end{textblock*}}

% Commands useful to place columns in the slides
\newcommand{\columnsbegin}{\begin{columns}}
\newcommand{\columnsend}{\end{columns}}
\newcommand{\columnhalf}{\column{.5\textwidth}}
\newcommand{\columnlarge}{\column{.66\textwidth}}
\newcommand{\columnsmall}{\column{.33\textwidth}}
\newcommand{\columnverylarge}{\column{.75\textwidth}}
\newcommand{\columnverysmall}{\column{.25\textwidth}}

% Copyright notices and sources:
% The background picture: http://www.uberpiglet.com/vectors/20-amazing-vector-backgrounds/
% UMONS-EcoNum template is partly inspired from a work done by:
% C. Troestler <Christophe.Troestler@umons.ac.be> licensed under GNU GPL v3 or later.

% SLIDE NOTES:
% For slide notes, use the \note{} beamer style in your document.
% To make sure to generate note pages even where there are no notes
% (in order to have side-by-side slides at left and notes at right),
% we add this:
\makeatletter
\def\beamer@framenotesbegin{%
  \gdef\beamer@noteitems{}%
  \gdef\beamer@notes{{}}%
}
\makeatother
% Uncomment this line to print notes:
%\setbeameroption{show notes}
% To interleave slides and notes in the PDF:
%\setbeamertemplate{note page}[plain]
% To get only notes:
%\setbeameroption{show only notes}

% Uncomment the following 3 lines in presentation mode for dual screen
% use (and specify 'left', 'top', 'right' or 'bottom' for second screen...)
%\usepackage{pgfpages}
%\setbeameroption{second mode text on second screen=right}
%\setbeameroption{show notes on second screen}

% More packages and options
\usepackage{graphicx}
\usepackage{rotating}
%\setbeamertemplate{caption}[numbered]
\usepackage{hyperref}
\usepackage[normalem]{ulem}
%\mode<presentation>
\usepackage{wasysym}
%\usepackage{amsmath}

% Redefine colors for links
\hypersetup{colorlinks=true,urlcolor=umons-red,linkcolor=umons-red,% link color controls section, subsection, and title
citecolor = umons-red, anchorcolor = umons-red}
% We do not want to change link colors in infolines
\addtobeamertemplate{headline}{\hypersetup{allcolors=.}}{}
\addtobeamertemplate{footline}{\hypersetup{allcolors=.}}{}

% Do not print slides with first level header
\AtBeginSection[]{}

% Redefine how subsection slides appear
\AtBeginSubsection{
  \frame{
    \centering
    \LARGE
    \insertsubsection
  }
}
\usepackage{bookmark}
\IfFileExists{xurl.sty}{\usepackage{xurl}}{} % add URL line breaks if available
\urlstyle{same}
% fallback for those not using the hyperref driver hyperxmp:
\makeatletter
\@ifundefined{xmpquote}{\newcommand{\xmpquote}[1]{#1}}{}
\makeatother
\hypersetup{
  pdftitle={Science des données II},
  pdfauthor={Philippe Grosjean \& Guyliann Engels},
  hidelinks,
  pdfcreator={LaTeX via pandoc}}

\title{\texorpdfstring{Science des données II}{Science des données II}}
\subtitle{\texorpdfstring{\includegraphics[width=.08\textwidth,height=.1\textheight]{../../template/biodatascience.png}
\includegraphics[width=.08\textwidth,height=.1\textheight]{../../template/SciViews-logo.pdf}
\vfill Reprise et installation du
logiciel}{  Reprise et installation du logiciel}}
\author{\texorpdfstring{Philippe Grosjean \& Guyliann
Engels}{Philippe Grosjean \& Guyliann Engels}}
\date{}
\institute{Université de Mons, Belgique\break Service d'Écologie
numérique\break \includegraphics[width=.08\textwidth,height=.1\textheight]{../../template/EcoNum-logo.pdf}
\break \url{https://wp.sciviews.org} \break \url{sdd@sciviews.org}}

\begin{document}
\frame{\titlepage}

\begin{frame}{Reprise\ldots{}}
\protect\phantomsection\label{reprise}
C'est la reprise\ldots{} pas encore tout-à-fait~! Durant ces deux
heures, nous allons essentiellement nous \textbf{préparer} au cours de
SDD II.

\center

\includegraphics[width=0.6\linewidth,height=\textheight,keepaspectratio]{B00Pa_intro_files/figures/start.jpg}

\begin{itemize}
\item
  Installation du logiciel
\item
  Rappel sur l'organisation des cours de SDD
\item
  Travail pour le prochain cours et nouveautés
\end{itemize}
\end{frame}

\begin{frame}{Installation et configuration des logiciels}
\protect\phantomsection\label{installation-et-configuration-des-logiciels}
\center

\includegraphics[width=0.4\linewidth,height=\textheight,keepaspectratio]{B00Pa_intro_files/figures/saturncloud.png}

Cette année, nous utiliserons la \textbf{SciViews Box 2026} dans le
Cloud avec \textbf{Saturn Cloud}.
\end{frame}

\begin{frame}{Étapes de l'installation}
\protect\phantomsection\label{uxe9tapes-de-linstallation}
\begin{enumerate}
\setcounter{enumi}{-1}
\tightlist
\item
  Accéder au cours sur Moodle \textbf{Science des données II :
  modélisation} et se connecter à son compte GitHub
  \url{https://github.com/}.
\item
  Aller sur le site du cours : \url{https://wp.sciviews.org/}
\item
  Lier le cours et se connecter avec son compte GitHub.
\item
  Rejoindre l'organisation Econum de Saturn Cloud (invitation par mail).
\end{enumerate}

Cliquer sur le bouton bleu \textbf{RStudio} en haut à droite puis suivez
les instructions d'installation.

\begin{enumerate}
\setcounter{enumi}{3}
\tightlist
\item
  Créer votre ressource : \textbf{svbox2026sdd2} (attention à son nom).
\end{enumerate}
\end{frame}

\begin{frame}{Gestion des séances : Rappel}
\protect\phantomsection\label{gestion-des-suxe9ances-rappel}
\begin{itemize}
\item
  \textbf{le plan de cours} est l'élément indispensable à connaitre
  (voir
  \href{https://github.com/BioDataScience-Course/BioDataScience-Common/tree/2026-2027/docs/plan_de_cours/sdd2_plan_cours_2026-2027.pdf}{plan
  de cours}).
\item
  \textbf{séances assimilées à des TP}. Obligatoires. Cela ne change
  pas.
\item
  \textbf{planning du travail} dans le cours (voir section
  \href{https://wp.sciviews.org/sdd-umons2/?iframe=wp.sciviews.org/sdd-umons2-2026/planning.html}{planning
  des séances}).
\end{itemize}
\end{frame}

\begin{frame}{Gestion des modules : Rappel}
\protect\phantomsection\label{gestion-des-modules-rappel}
\begin{itemize}
\item
  Chaque module est travaillé pendant \textbf{une semaine}
\item
  \textbf{Deux séances de 2h et de 4h} en présentiel, respectivement et
  12 à 15h de travail au total pour l'étudiant (= 0.5 ECTS) par module
\item
  Un module \textbf{une semaine sur deux} pour laisser le temps de
  finaliser le précédent et de préparer le suivant
\end{itemize}

\includegraphics[width=0.8\linewidth,height=\textheight,keepaspectratio]{B00Pa_intro_files/figures/semaine_type.pdf}
\end{frame}

\begin{frame}{Évaluation continue : Rappel}
\protect\phantomsection\label{uxe9valuation-continue-rappel}
\textbf{Évaluation continue} prenant en compte \textbf{toute votre
activité} (\emph{chaque} exercice donne des points)~:

\begin{itemize}
\tightlist
\item
  Exercices N1 \& 2 en distanciel
\item
  Projets individuels et en groupe
\item
  Une évaluations sommative par module
\end{itemize}

\pandocbounded{\includegraphics[keepaspectratio]{B00Pa_intro_files/B00Pa_intro_files/figure-beamer/unnamed-chunk-1-1.pdf}}
\end{frame}

\begin{frame}{Préparation pour le prochain cours}
\protect\phantomsection\label{pruxe9paration-pour-le-prochain-cours}
\center

\includegraphics[width=0.4\linewidth,height=\textheight,keepaspectratio]{B00Pa_intro_files/figures/notes.jpg}

\begin{itemize}
\item
  Revoir ce qui n'était pas compris ou oublié de SDD I
\item
  Se mettre à jour pour ceux qui ne le sont pas
\item
  Exceptionnellement, il ne faut \emph{pas} préparer le module 1. Il est
  assez léger et nous le travaillerons ensemble en séance (et
  actuellement, il n'est pas encore complètement finalisé).
\end{itemize}

\textbf{Ressources utiles~:}

\begin{itemize}
\tightlist
\item
  Site web du cours : \url{http://bds.sciviews.org/}
\item
  Cette présentation~:
  \url{https://github.com/BioDataScience-Course/sdd_lessons/tree/2026-2027/B00/presentations}
\end{itemize}
\end{frame}

\begin{frame}{Nouveauté : Tutoriels}
\protect\phantomsection\label{nouveautuxe9-tutoriels}
\begin{itemize}
\item
  Nouvelle SciViews Box 2026 comme workspace dans une interface remaniée
  sur Saturn Cloud
\item
  GitHub Classroom n'existant plus, nous créons nous-mêmes vos dépôts et
  vous aurez uniquement un lien URL vers GitHub pour y accéder.
\end{itemize}
\end{frame}

\end{document}
